\begin{trapbox}

	\begin{description}[style=nextline, leftmargin=2.8em, labelsep=0.5em, itemsep=0.35em]
		\item[\textbf{\textcolor{CTrap}{易错点一（基底条件混淆）}}]
			在非正交基底中，仍套用 $x_1x_2+y_1y_2$ 计算点乘.
		\item[\textbf{\textcolor{CTrap}{正确理解（正交基底限定）：}}]
			点乘坐标公式只在单位正交基底（直角坐标系）下成立.
		\item[\textbf{\textcolor{CTrap}{错因分析（忽视前提）：}}]
			公式推导依赖了 $\mathbf{i}\cdot\mathbf{j}=0$，忽略这个前提会导致错误.
	\end{description}

	\medskip
	\begin{description}[style=nextline, leftmargin=2.8em, labelsep=0.5em, itemsep=0.35em]
		\item[\textbf{\textcolor{CTrap}{易错点二（符号顺序颠倒）}}]
			计算 $\vec{a}-\vec{b}$ 时，误写成 $x_2-x_1$ 或 $y_2-y_1$.
		\item[\textbf{\textcolor{CTrap}{正确理解（按序对应相减）：}}]
			应为 $x_1-x_2$ 和 $y_1-y_2$，注意是被减向量坐标在前.
		\item[\textbf{\textcolor{CTrap}{错因分析（惯性思维颠倒）：}}]
			未注意减法没有交换律，直接套用加法模式.
	\end{description}

	\medskip
	\begin{description}[style=nextline, leftmargin=2.8em, labelsep=0.5em, itemsep=0.35em]
		\item[\textbf{\textcolor{CTrap}{易错点三（点乘结果属性）}}]
			把 $\vec{a}\cdot\vec{b}$ 当作向量，继续与其他向量进行点乘或数乘时混淆.
		\item[\textbf{\textcolor{CTrap}{正确理解（实数非向量）：}}]
			$\vec{a}\cdot\vec{b}$ 是一个实数，不能再作点乘对象（与数乘不同）.
		\item[\textbf{\textcolor{CTrap}{错因分析（混淆点乘与数乘）：}}]
			对数乘和点乘的运算对象没分清：数乘是实数乘向量，点乘是向量乘向量得实数.
	\end{description}

\end{trapbox}
